\section{Lecture 7 (September 25th)}
\begin{thm}
(Schwartz-Pick theorem) Let $f:U\rightarrow U$ be an analytic function such that $f(\alpha )=\beta $. Then,
\[|f'(\alpha )|\leq \dfrac{1-|\beta |^2}{1-|\alpha |}\]
and the equality holds when 
\[f(z)=\phi _{-\beta }(\lambda \phi _{\alpha }(z))\]
for some $\lambda \in T$. 
\end{thm}
\vspace{2ex}
\begin{ex}
Observe that $f(x)=x^3$ is one-to-one on ${\bm R}$, but $f'(0)=0$. Therefore, we see that on ${\bm R}$, $f$ being one-to-one on an interval $I$ doesn't necessarily imply that $f'(x)=0$ for any $x\in I$. However, the converse works as if $f'(x)\ne 0$ for any $x\in I$, then $f$ is one-to-one by the mean value theorem.

\bigksip

However, on ${\bm C}$, $f(z)=e^{z}$ satisfies $f'(z)\ne 0$ for any $z\in {\bm C}$ but $f$ is not one-to-one. We claim that if $f$ is analytic and one-to-one on a domain $D$, then $f'(z)\ne 0$ for any $z\in D$.
\end{ex}
\vspace{2ex}
\begin{thm}
If $f:U\rightarrow U$ is one-to-on analytic, meaning that $f'(z)\ne 0$ for any $z\in U$, then $f(z)=e^{i\theta }\phi _{\alpha }(z)$ for some $\alpha \in U$, $\theta \in {\bm R}$
\end{thm}
\vspace{2ex}
\begin{proof}
Let $f(\alpha )=0$. First define $g=f^{-1}$ on $U$. Then $g(f(z))=z$ for all $z\in U$ and $g$ is analytic on $U$. Due to the fact that the derivative is never zero, we can use the inverse function theorem, where we are able to write
\[\dfrac{g(w)-g(w_0)}{w-w_0}=\dfrac{z-z_0}{f(z)-f(z_0)}=\dfrac{1}{\dfrac{f(z)-f(z_0)}{z-z_0}}\]
and
\[g'(w_0)=\dfrac{1}{f'(z_0)}\]
since $f'(z_0)\ne 0$. We see that
\[g'(0)f'(\alpha )=1\]
by the chain rule. By the Schwartz-Pick theorem,
\[|f'(\alpha )|\leq \dfrac{1}{1-|\alpha |^2}\]
and $|g'(0)|\leq 1-|\alpha |^2$. This means that 
\[\dfrac{1}{1-|\alpha |^2}\leq |f'(\alpha )|\leq \dfrac{1}{1-|\alpha |^2}\]
By the Schwartz lemma,
\[f(z)=\lambda \phi _{\alpha }(z)\]
for some $\lambda \in T$. 
\end{proof}
\vspace{2ex}
\begin{thm}
If $f$ is analytic in a domain $D$ with $f'(z_0)\ne 0$ for some $z_0\in D$ then there is $r>0$ such that $D(z_0,r)\subset D$ and $f$ is one-to-one on $D(z_0,r)$. 
\end{thm}
\vspace{2ex}
\begin{proof}
If $z,w\in D(z_0,r)$ and $z\ne w$, then 
\[\Big|\dfrac{f(z)-f(w)}{z-w}-f'(z_0)\Big|<\dfrac{1}{2}|f'(z_0)|\]
which implies that
\[\dfrac{1}{2}|f'(z_0)|<\Big|\dfrac{f(z)-f(w)}{z-w}\Big|\]
if $z,w\in D(z_0,r)$ with $z\ne w$. 
\end{proof}
\vspace{2ex}
\begin{lem}
If $f$ is analytic in a domain $D$ and 
\[g(z,w)=\begin{cases}
\dfrac{f(z)-f(w)}{z-w}&z\ne w\\
f'(z)&z=w
\end{cases}\]
then $g$ is continuous on $D\times D$. Note that this is trivially true on ${\bm R}$.
\end{lem}
\vspace{2ex}
\begin{lem}
If $|b-a|<\dfrac{|a|}{2}$, then $|b|> \dfrac{|a|}{2}$.
\end{lem}
\vspace{2ex}
\begin{proof}
Notice that if suffices to prove the statement only when $z=w$. Since $f'$ is continuous on $D$, for every $\varepsilon $, there exists a $\delta $ such that if $z\in D(a,\delta )$ then $|f'(z)-f'(a)|<\varepsilon $. Take any $z,w$ ($z\ne w$) in $D(a,r)$. Define $\xi (t)=tz+(1-t)w$ for $0\leq t\leq 1$. Then,
\[\int ^{1}_{0}f'(\xi (t))\,dt=\dfrac{f(z)-f(w)}{z-w}=g(z,w)\]
so that
\[g(z,w)-g(a,a)=\int ^{1}_{0}[f'(\xi (t))-f'(a)]\,dt\]
Since $z,w\in D(a,r)$ we have
\[\Big|\int ^{1}_{0}[f'(\xi (t))-f'(a)]\,dt\Big|\leq \int ^{1}_{0}|f'(\xi (t))-f'(a)|\,dt< \varepsilon \]
In conclusion, if $z,w\in D(a,r)$, then $|g(z,w)-g(a,a)|<\varepsilon $. 
\end{proof}
\vspace{2ex}
\begin{thm}
(Rouche's theorem) If $f$ and $g$ are analytic inside and on a simple, closed contour $C$ such that $|g(z)|<|f(z)|$ for every $z\in C$, then $f$ and $f+g$ have the same number of zeros inside $C$.
\end{thm}
\vspace{2ex}
\begin{thm}
(Open mapping theorem) Let $f$ be non-constant analytic in a domain $\mathrm{\Omega}$. If $D\subset \mathrm{\Omega} $ is open, then $f(D)$ is open.
\end{thm}
\vspace{2ex}
\begin{proof}
To show that $f(D)$ is open, take any $w_0\in f(D)$. We need to find some $m>0$ such that $D(w_0,m)\subset f(D)$. Let $f(z_0)=w_0$. Then there is $r>0$ such that
\[\bar{D}(z_0,r)\subset D\]
and $f(z)-w_0$ has no zeros in $0<|z-z_0|\leq r$ by the identity theorem. Let
\[m=\min _{|z-z_0|=r }|f(z)-w_0|>0\]
Then we'll show that $D(w_0,m)\subset f(D)$. Take any $w_1\in D(w_0,m)$ (we're trying to show that $w_1\in f(D)$). Then
\[|w_0-w_1|<|f(z)-w_0|\]
on $|z-z_0|=r$. Rouche's theorem implies that $f(z)-w_0$ and $f(z)-w_1$ have the same number of zeros inside $|z-z_0|=r$ which implise that
\[w_1\in f(D(z_0,r))\subset f(D)\]
\end{proof}
\vspace{2ex}
\begin{thm}
If $f$ is one-to-one analytic in a domain $D$, then $f'(z)\ne 0$ for any $z\in D$. 
\end{thm}
\vspace{2ex}
\begin{proof}
Suppose that $f'(z_0)=0$ for some $z_0\in D$ (we'll show that $f$ is not one-to-one). Then $f(z)-f(z_0)=f''(z_0)(z-z_0)^2+\ldots $ in some neighborhood of $z_0$. There is $r>0$ such that $\bar{D}(z_0,r)\subset D$ and $f(z)-f(z_0)$ and $f'(z)$ both have no zeros in $0<|z-z_0|\leq r$.

\bigskip

Let $m=\min_{|z-z_0|=r}|f(z)-f(z_0)|$ and choose any $w_1\in D(f(z_0),m)$. Then,
\[|f(z_0)-w_1|<m\leq |f(z)-f(z_0)|\]
on $|z-z_0|=r$. By Rouche's theorem, $f(z)-f(z_0)$ and $f(z)-w_1$ have the same number of zeros in $D(z_0,r)$ where $f(z)-f(z_0)$ has at least 2 zeros. In other words, since $f'\ne 0$ in $D(z_0,r)$, there is at least 2 distinct zeros of $f(z)-w_1$ in $D(z_0,r)$ where $f$ is not one-to-one.
\end{proof}
\vspace{2ex}

